\section{Probability}
\label{sec:theory_general_probability}

Probability is your best friend when dealing with real-world
electronic components. All electronic device properties suffer
from variations. Namely there are two types of variations:
\begin{enumerate}
  \item \textbf{Systematic Variations:} These are predictable
    variations that occur due to known factors in the
    manufacturing process. For example, if a certain batch of
    resistors is consistently \(2\%\) higher in resistance due
    to a calibration error in the manufacturing equipment, this
    is a systematic variation.
  \item \textbf{Random Variations:} These are unpredictable
    variations that occur due to inherent randomness in the
    manufacturing process. For example, slight differences in
    material properties or environmental conditions during
    manufacturing can lead to random variations in component
    values.
\end{enumerate}
It is almost impossible to elimiate variations if you're not making
the components yourself, but as a designer you can take steps to
understand and mitigate their impact on your designs.

\subsection{Pre and Post Production}
Before production, designers use statistical models like Monte Carlo
simulations to predict how these variations will affect circuit
performance. After production, manufacturers perform quality control
tests to ensure that components meet specified tolerances.

In both cases, It is done by analysing a large number of components,
recording the values for interesting parameters, and analyzing the
resulting distribution.

\subsection{Available Tools}
It is common to assume that these variations follow a normal (Gaussian)
distribution, although this may not always be the case. There are several
tools available to analyze and interpret these distributions, including:
\begin{itemize}
  \item \textbf{Mean (Average):} The mean is the average value
    of all measured component values. $$\mu = \frac{1}{N} \sum_{i=1}^{N} x_i$$
    where \(N\) is the total number of measurements and \(x_i\)
    represents each individual measurement.
  \item \textbf{Variance (\(\sigma^2\)):} Variance measures how
    much the component values deviate from the mean. It is
    calculated as the average of the squared differences from
    the mean. $$\sigma^2 = \frac{1}{N} \sum_{i=1}^{N} (x_i - \mu)^2$$
  \item \textbf{Standard Deviation (SD or \(\sigma\)):} The
    standard deviation is the square root of the variance and
    provides a measure of the spread of component values around
    the mean. $$\sigma = \sqrt{\sigma^2}$$
  \item \textbf{Cumulative Distribution Function (CDF):} The CDF
    represents the probability that a random variable takes on
    a value less than or equal to a specific value. It is useful
    for determining the likelihood of a component falling within
    a certain range.
  \item \textbf{Probability Density Function (PDF):} The PDF
    describes the relative likelihood of a random variable to
    take on a specific value. It is the derivative of the CDF and
    provides insights into the distribution of component values.
\end{itemize}

\subsection{Mitigations}
For an analog integrated circuit designer, these tricks can help
mitigate the impact of variations:
\begin{itemize}
  \item \textbf{Redundancy:} Use components with fused
    nets that can be trimmed or adjusted after manufacturing to
    achieve the desired performance.
  \item \textbf{Common Centroid Layout:} A common centroid layout
    technique can help reduce the impact of systematic variations
    by averaging out the effects across multiple components.
  \item \textbf{Area Scaling:} M. Pelgrom discovered that many
    variations decrease as the area of the component increases.
    This is known as Pelgrom's Law. By increasing the size of
    components, you can reduce the impact of random variations
    (by approximately the square root of the area).
  \item \textbf{Dummy Structures:} Adding dummy structures around
    critical components can help create a more uniform
    manufacturing environment, reducing variations caused by
    edge effects.
\end{itemize}
